\scanpage{25}
\mspage{\handwritten{\unsure{11|15}\deleted{\unsure{a}}}}
\deleted{\unsure{G de R.|de R.} En vertu de 10.5.3. c'est un sous-groupe fermé de $G$ lisse sur $S$, donc lisse et affine sur $S$, et comme $T$ est central, on a $T\unsure{\subset|=}H$.}

Ceci résulte du \deleted{résultat suivant}

\uline{Lemme} \handwritten{1}.14. Soient $G$ \deleted{affine}\inserted{\handwritten{un $S$-préschéma en groupes,}} de présentation finie sur $S$, $T$ un tore maximal de $G$, \deleted{alors $T$ contient tout sous-tore \uline{central}} $R$ un sous-tore de $G$, et supposons que $T$ et $R$ commutent. Alors $R\handwritten{\subset} T$.

En effet, comme $R$ et $T$ commutent, le morphisme $R\times_S T\handwritten{\to}G$ défini \deleted{\unsure{pa|par}} par $(r,t)\handwritten{\rightsquigarrow}r.t$ est un homomorphisme de groupes, donc en vertu de \handwritten{1}.8. il admet un sous-groupe image $T'$ dans $G$, qui est un \deleted{sous-}groupe de type multiplica\unsure{ti|it}f quotient de $R\times_S T$, donc un tore, contenant évidemment $T$. Comme $T$ est un tore maximal, on a $T'=T$ (\handwritten{1}.4.), donc $R\handwritten{\subset}T$, \deleted{cqfd.} ce qui démontre \handwritten{1}.14. donc \handwritten{1}.13. En particulier\handwritten{, utilisant 6.7. b):}

\uline{Corollaire} \handwritten{1}.15. Soit $G$ un $S$-préschéma en groupes \deleted{\unsure{affi|affine}} \uline{commutatif}, lisse et affine sur $S$\handwritten{,} \inserted{\handwritten{de rang réductif localement constant.}} Alors $G$ admet u\deleted{\unsure{j}}n unique tore maximal, et c\deleted{\unsure{l}}e\deleted{\unsure{lui|t}} \deleted{qui est le plus} dernier contient tout sous-tore de $G$.

\uline{Corollaire} \handwritten{1}.16. Soit $G$ un $S$-préschéma en groupes, lisse et affine sur $S$. Supposons que pour tout $s\handwritten{\in}S$, désignant par $\bar s$ le spectre d'une clôture algébrique \inserted{$k$} de $k(s)$, la fibre géométrique $G_{\bar s}$ soit un groupe algébrique connexe, et nilpotent (au sens de BIBLE, i.e. le groupe $G(k)$ de ses points à va\deleted{\unsure{b}}leurs dans $k$ est nilpotent). \inserted{Supposons de plus le rang réductif de $G$ loc.constant.} Alors $G$ admet u\deleted{\unsure{j}}n unique tore maximal $T$, \inserted{\handwritten{de plus}} $T$ est central et c'est le plus grand sous-tore de $G$.

En vertu de \handwritten{1}.7. b), $G$ admet localement pour fpqc un tore maximal, et en vertu de \handwritten{1}.13. on est ramené à prouver qu'un tore maximal de $G$ est central. En vertu de Exp \handwritten{$\overline{\underline{\mathrm{IX}}}$}\deleted{\illegible{1}}, 5.6. b), on est ramené au cas où $S$ est le spectre d'un corps, qu'on peut supposer algébriquement clos.
% archive sheet number: 24
